% Options for packages loaded elsewhere
\PassOptionsToPackage{unicode}{hyperref}
\PassOptionsToPackage{hyphens}{url}
%
\documentclass[
  12pt,
]{article}
\usepackage{amsmath,amssymb}
\usepackage{iftex}
\ifPDFTeX
  \usepackage[T1]{fontenc}
  \usepackage[utf8]{inputenc}
  \usepackage{textcomp} % provide euro and other symbols
\else % if luatex or xetex
  \usepackage{unicode-math} % this also loads fontspec
  \defaultfontfeatures{Scale=MatchLowercase}
  \defaultfontfeatures[\rmfamily]{Ligatures=TeX,Scale=1}
\fi
\usepackage{lmodern}
\ifPDFTeX\else
  % xetex/luatex font selection
\fi
% Use upquote if available, for straight quotes in verbatim environments
\IfFileExists{upquote.sty}{\usepackage{upquote}}{}
\IfFileExists{microtype.sty}{% use microtype if available
  \usepackage[]{microtype}
  \UseMicrotypeSet[protrusion]{basicmath} % disable protrusion for tt fonts
}{}
\makeatletter
\@ifundefined{KOMAClassName}{% if non-KOMA class
  \IfFileExists{parskip.sty}{%
    \usepackage{parskip}
  }{% else
    \setlength{\parindent}{0pt}
    \setlength{\parskip}{6pt plus 2pt minus 1pt}}
}{% if KOMA class
  \KOMAoptions{parskip=half}}
\makeatother
\usepackage{xcolor}
\usepackage[margin=1in]{geometry}
\usepackage{graphicx}
\makeatletter
\def\maxwidth{\ifdim\Gin@nat@width>\linewidth\linewidth\else\Gin@nat@width\fi}
\def\maxheight{\ifdim\Gin@nat@height>\textheight\textheight\else\Gin@nat@height\fi}
\makeatother
% Scale images if necessary, so that they will not overflow the page
% margins by default, and it is still possible to overwrite the defaults
% using explicit options in \includegraphics[width, height, ...]{}
\setkeys{Gin}{width=\maxwidth,height=\maxheight,keepaspectratio}
% Set default figure placement to htbp
\makeatletter
\def\fps@figure{htbp}
\makeatother
\setlength{\emergencystretch}{3em} % prevent overfull lines
\providecommand{\tightlist}{%
  \setlength{\itemsep}{0pt}\setlength{\parskip}{0pt}}
\setcounter{secnumdepth}{-\maxdimen} % remove section numbering
\ifLuaTeX
  \usepackage{selnolig}  % disable illegal ligatures
\fi
\usepackage{bookmark}
\IfFileExists{xurl.sty}{\usepackage{xurl}}{} % add URL line breaks if available
\urlstyle{same}
\hypersetup{
  pdftitle={Professional resume of Stephen D. Simon},
  hidelinks,
  pdfcreator={LaTeX via pandoc}}

\title{Professional resume of Stephen D. Simon}
\author{}
\date{\vspace{-2.5em}}

\begin{document}
\maketitle

\subsection{Contact information}\label{contact-information}

\begin{itemize}
\tightlist
\item
  14814 Granada Ct, Leawood, KS 66224
\item
  913-814-9942 (home)
\item
  913-912-2076 (cell)
\item
  \href{mailto:mail@pmean.com}{\nolinkurl{mail@pmean.com}}
\end{itemize}

\subsection{Personal}\label{personal}

\begin{itemize}
\tightlist
\item
  Born June 22, 1956. Married with one child.
\end{itemize}

\subsection{Education}\label{education}

\begin{itemize}
\tightlist
\item
  Ph.D., University of Iowa, Statistics, December 1982
\item
  M.S., University of Iowa, Statistics, December 1978
\item
  B.A., University of Iowa, Mathematical Sciences, July 1977
\end{itemize}

\subsection{Experience (in reverse chronological
order)}\label{experience-in-reverse-chronological-order}

\begin{itemize}
\item
  Full Professor, Department of Biomedical and Health Informatics,
  University of Missouri-Kansas City (UMKC), February 2009 to present.
  Note: I was part-time until January 2022, when I became full time. In
  this position, I started the Research and Statistical Consult Service.
  With a second consultant, we helped with the planning of research
  studies, analyzed the data coming from these studies, and co-authored
  research presentations and publications. We helped researchers obtain
  internal and external grant funding. I left work with the Research and
  Statistical Consult Service in January 2016 because of the demands of
  a major research grant, but the service is still being run by another
  faculty member. In addition to this work, I helped with the
  introduction of a new Master's degree in Bioinformatics and a graduate
  certificate in clinical research. I helped develop three new courses
  (Introduction to R, Introduction to SAS, and Introduction to SPSS)
  including the preparation of materials for an online version of these
  classes. I have provided guest lectures for other classes (Responsible
  Conduct of Research, Clinical Trials).
\item
  Independent Statistical Consultant, P.Mean Consulting, July 2008 to
  present. In this position, I have provided consulting on research
  design and data analysis to clients. I have specialized in helping
  doctoral students who are struggling with the requirements of their
  dissertations. I have provided webinar classes for The Analysis
  Factor, Novartis, and the Medical Library Association. When I first
  started my consulting business, I could not find much information
  about the practical aspects of statistical consulting (e.g.,
  contracts, insurance, and billing), so I developed resources on my
  website and blog and gave presentations and short courses on setting
  up a consulting practice.
\item
  Research Biostatistician, Office of Medical Research, The Children's
  Mercy Hospital and Clinics (CMH), May 1996 to October 2008, with a
  joint appointment as Associate Professor (promoted to Full Professor,
  September 2007) in the School of Medicine at the University of
  Missouri-Kansas City. In this position, I provided statistical
  consulting to the doctors, nurses, and other health care professionals
  at CMH. Before I arrived, there was no statistical support for
  researchers, so I developed a consulting system from scratch. This
  included a series of training classes that helped researchers to
  become self-sufficient for some of their less complex data analyses. I
  wrote some informational handouts for clients and placed these on my
  web (now at www.pmean.com). It started slowly, but averaging one or
  two handouts per week, I now have a website that has over 2,000 pages.
  This website is regularly cited in both paper and electronic resources
  for research methodology and pops up high on the list of search
  engines when technical terms in research are entered. I helped many
  researchers get funding (both internal and external) through research
  grants. I worked closely with the Institutional Review Board to ensure
  that protocols submitted to them addressed concerns about scientific
  merit and protection of confidentiality. I provided training in the
  ethical conduct of research. I worked closely with the effort to
  implement Evidence Based Health Care at CMH and participated in
  numerous journal clubs. In 2006, I published a book on the statistical
  issues associated with critical appraisal of research evidence. I was
  the sole statistician at CMH from 1996 to 2007 and oversaw a massive
  expansion in the quantity and quality of research conducted there.
\item
  Chief, Statistics Activity, Division of Biomedical and Behavioral
  Science, National Institute for Occupational Safety and Health,
  December 1987 to May 1996. I supervised a staff of two research
  statisticians and two computer specialists. I oversaw the
  administration of a Local Area Network for more than 100 users, and
  the maintenance of administrative and research data on a client-server
  database. I reviewed and approved all project protocols in the
  Division prior to the start of any experimental work. I worked
  directly with project officers in designing experiments, analyzing
  data, and writing results for publication. I was very active in the
  NIOSH effort to adopt Total Quality Management (TQM), serving as team
  leader, facilitator, or member of a variety of quality teams. I
  provided many informal training seminars in TQM for NIOSH employees.
\item
  Research Statistician, Division of Biomedical and Behavioral Science,
  National Institute for Occupational Safety and Health, July 1987 to
  November 1987. The person who hired me for this position left shortly
  after I arrived and I was quickly promoted to her position (see
  above).
\item
  Assistant Professor, Department of Applied Statistics and Operations
  Research, Bowling Green State University, September 1981 to May 1987.
  Duties included teaching and research. I taught Business Statistics,
  Experimental Design, Linear Models, Multivariate Statistics,
  Nonparametric Statistics, and Linear Regression. I published articles
  and made presentations in the areas of numerical accuracy, robust
  statistics, and computer applications. During the Summer of 1986, I
  served as the Assistant Director of the Statistical Consulting Center.
\item
  Adjunct Instructor, Department of Statistics, University of Iowa, June
  to August 1981. I taught a three credit hour class, Statistical
  Computing.
\item
  Research Assistant, Statistical Consulting Center, University of Iowa,
  June 1979 to August 1981. I was the student director of the
  Statistical Consulting Center. Under the direction of a faculty
  adviser, I met with faculty and graduate students and provided
  guidance and support for the design of research studies and the
  analysis of the data from those studies.
\end{itemize}

\subsection{List of research grants}\label{list-of-research-grants}

{[}19{]} Public Health and the Built Environment (Housing and Urban
Development, Principal Investigator, Neal Wilson) I am a co-investigator
funded at 10\% effort from January 2025 to December 2028.

{[}18{]} Faithful Response II: COVID-19 Rapid Test-to-Treat with African
American Churches (U01MD018310, Principal Investigator: Jannette
Berkley-Patton) The primary aim of this study is to fully test a
culturally/religiously-tailored, church-based COVID-19 test-to-treat
with linkage-to-care intervention condition against a non-tailored
education condition on COVID-19 rapid testing rates at 6 months with
adult African-American church members and the community members they
serve. I am a co-investigator funded at 10\% effort from November 2022
to October 2024.

{[}17{]} Impact Lead - Kansas City (HHTS20000276, Housing and Urban
Development, Principal Investigator: Steve Simon). We propose a data
collection effort to show that lead safe interventions can be targeted
more effectively and that they can inform decisions about the choices
among intervention strategies. Our first aim is to quantify the extent
to which remediations to make housing lead safe by HUD standards leads
to fewer lead poisoned children among those who move in. Our second aim
is to develop an exterior housing-based risk index to cost effectively
target homes with higher interior lead dust levels. I am the principal
investigator funded at 13.5\% effort from January 2021 through December
2024.

{[}16{]} Development of a telehealth obesity intervention for patients
with MS. (Multiple Sclerosis Society, Principal Investigator: Jared
Bruce) A cross-disciplinary research team is testing the effectiveness
of an MS-specific weigh loss/healthy living program delivered by phone,
since obesity can profoundly worsen MS severity. I am a co-investigator
funded at 1.5\% effort starting in October 2019, increasing to 5\%
effort from October 2022 through September 2023 and dropping back to
1.5\% effort from October 2023 to September 2024.

{[}15{]} Our Healthy Kansas City Eastside (Jackson County funding,
Principal Investigator: Jannette Berkley-Patton). This is a
community-wide initiative that promotes and delivers widespread COVID-19
vaccinations and other health services to residents on the east side of
Kansas City. More than 60 community organizations and health agencies
are partnering with us to support healthy lifestyles through vaccine
events and health screenings like blood pressure checks, diabetes
screenings and dental education. I was a co-investigator funded at 5\%
effort from July 2022 through March 2023.

{[}14{]} COVID-19 Testing and Linkage to Care with African-American
Church and Health Agency Partners (DK124664-01S1, National Institute of
Diabetes and Digestive and Kidney Diseases, Principal Investigator:
Jannette Berkley-Patton). The primary aim of this study is to fully test
a culturally/religiously-tailored, church-based COVID19 testing and
linkage to care (LTC) intervention condition against a non-tailored
intervention condition on COVID19 testing rates at 6 months with adult
AA church members and the community members they serve. Guided by the
Theory of Planned Behavior and Socioecological Model, our community-
engaged approach includes trained church leaders delivering a
culturally, church-appropriate COVID19 Toolkit. Examination of LTC use
and contact tracing approval will aid in understanding intervention
impact on COVID19 testing by addressing participant essential needs.
Potential mediators/moderators related to receipt of COVID-19 testing
will be evaluated, and a process evaluation to determine implementation
facilitators, barriers, and fidelity related to increasing COVID19
testing rates. I was a co-investigator funded at 10\% effort during
calendar year 2021 and 5\% effort during calendar year 2022.

{[}13{]} Preschool Development Grant Birth through Five Initiative
(Missouri Department of Elementary and Secondary Education, 90TP0066-01,
Principal Investigator: Mike Abel). We present best practices for
professionals and programs across Missouri's early childhood
mixed-delivery system. For more than 50 years, early childhood programs
that support the growth and development of young children have been
shown to have a profound impact on children over their lifespan. The
challenge is to identify what elements of a high-quality early childhood
program are most associated with positive outcomes for children and
families. I was a co-investigator funded at 5\% effort from September to
December 2020.

{[}12{]} Frontiers: University of Kansas Clinical and Translational
Science Institute (National Institutes of Health, 5UL1TR002366-04,
Principal Investigator: Mario Castro). The Heartland Institute for
Clinical and Translational Research has been a catalyst for bringing
together translational science investigators and stakeholders across the
KC region, and beyond. The vision of Frontiers is to contribute to and
lead national efforts to transform the way we do clinical and
translational research (CTR), and to ensure research is more rapidly and
more efficiently translated to the point of care so that it may
contribute to improved health. I was a co-investigator funded at 10\%
effort from September 2017 through June 2021.

{[}11{]} Greater Plains Collaborative (Patient Centered Ooutcomes
Research Institute, HSRP20162063, Principal Investigator: Lemuel
Waitman). The Greater Plains Collaborative (GPC) is a new network of 10
leading medical centers in 7 states committed to a shared vision of
improving healthcare delivery through ongoing learning, adoption of
evidence-based practices, and active research dissemination. Role:
Consultant. I billed at an hourly rate from February 2016 to August 2017
when my efforts were switched to ``Frontiers: University of Kansas
Clinical and Translational Science Institute'' described above.

{[}10{]} Increasing HIV Screening in African American Churches.
(National Institutes of Health, R01 MH099981, Principal Investigator:
Jannette Berkley-Patton). The primary aim of the proposed study is to
conduct a clustered, randomized community trial to fully test our
culturally/religiously-tailored church-based HIV screening intervention
against a standard information intervention on HIV screening rates at 6
and 12 months with adult AA church members and community members who use
church outreach services. Role: Co-investigator. I was a co-investigator
funded at 5\% effort from July 2013 through January 2016, but had to
discontinue work on this grant because of other work commitments.

{[}9{]} Role of cytochrome P450 in alcohol-mediated effects on
antiretrovirals and HIV-1. (National Institutes of Health, R01 K2304028,
Principal Investigator: Santosh Kumar). Our central hypothesis is that
in monocytes/macrophages alcohol-mediated oxidative stress causes
enhanced HIV-1 replication via pathways mediated through CYP2E1, and
alcohol-mediated decrease in the efficacy of PIs and increase in
PIs-mediated toxicity concurrent with increase in the rates of HIV-1
replication is directly mediated through CYP3A4. Role: Co-investigator.
I was a co-investigator funded at 8\% effort starting in September 2013
with increasing support in future years, but support was lost in July
2014 when the PI moved to a different institution.

{[}8{]} HIV Testing in African American Churches. (National Institutes
of Health, K01 MH082640, Principal Investigator: Janette
Berkley-Patton). This Mentored Research Career Award provides support
for Dr.~Jannette Berkley-Patton to obtain the requisite training and
expertise to develop and test a church-based intervention to increase
HIV testing among African Americans. Role: Co-investigator. I was a
co-investigator funded at 5\% effort from April 2012 through July 2013
when my efforts were switched to ``Increasing HIV Screening in African
American Churches'' described above.

{[}7{]} Ensuring Medication Safety in Pediatric Emergencies. (Food and
Drug Administration, U01 FD004249, Principal Investigator, Susan
Rahman). In no other patient population is there more calculation and
manipulation of drug doses than in children. The need to individualize
pediatric doses introduces the potential for errors along the entire
process of drug delivery from prescribing and transcribing the dose to
diluting, compounding, dispensing and administering the final
formulation. Consequently, dosing errors represent the most common type
of medication error in children. Error rates in children are 3x more
common than the corresponding rates in adults and in hospital settings
the rate of medication errors on pediatric wards is second only to those
in intensive care units. Despite the abundance of weight estimation
strategies, there remains a critical need for methods that are accurate
across a wide range of pediatric ages, weights, lengths, nationalities
and body compositions. Consequently, we undertook the development and
validation of a weight estimation method that incorporates surrogates
for both stature and body habitus, requires no subjective assessment and
performs robustly independently of age and length over a broad range of
weights. Role: Consultant I billed at an hourly rate from September 2011
through September 2012.

{[}6{]} Health Educators in the Emergency Department to Improve Sexual
Health Care for Adolescents. (National Institutes of Health, KL2
RR033177 \& UL1 RR033179, Principal Investigator: Melissa Miller). Our
objective in this application is to develop an ED-based intervention for
adolescents aged 14-19 years that will increase access to care, provide
more timely detection of STIs and pregnancy, and ultimately reduce
adolescents' risk for STIs and unintended pregnancies. Role: Consultant.
I billed at an hourly rate from August 2011 through July 2012.

{[}5{]} CAM, Medical Service Utilization, and Quality of Care among
Patients with Musculoskeletal Conditions (National Institutes of Health,
R01 AT000891, Principal Investigator: Bill Lafferty). To inform health
policy decision making at the individual and societal level by assessing
the effect of complementary and alternative medicine (CAM) provider use
on three primary outcomes: costs, service utilization, and quality of
care for the treatment of musculoskeletal conditions. I was a
co-investigator funded at 25\% effort from November 2010 to August 2011.

{[}4{]} NDNQI Services Agreement Work Order No.~3 (Principal
Investigator: Nancy Dunton). Funded by the American Nurses Association.
The major goals of this project are to develop a method for unit-based
acuity adjustment and to develop a method to calculate hospital level
indicators. I was a co-investigator funded at 25\% effort from February
through April 2008.

{[}3{]} Development of a Sperm Morphology Standard (National Institutes
of Health, R43 HD044383, Principal Investigator: Susan Rothmann). Total
direct costs of \$137,349 in 2003. Infertility affects an estimated ten
million Americans annually. A critical ``gateway'' diagnostic test used
in over 10,000 USA labs is sperm morphology, one of the most predictive
measures of fertility potential that often determines the course of
clinical intervention. Morphology also is an important measure of toxic
exposure from environmental and occupational sources. Unfortunately,
morphology analysis often is the most unreliable and difficult part of a
semen profile to interpret, due primarily to lack of a useful standard.
Several schemes exist, none of which have adequate agreement among
laboratories. The goal of the proposal is to determine the feasibility
of developing a consensus Standard for sperm morphology associated with
reproducible visual and morphometric characteristics. I was funded as a
consultant on an hourly basis from May 2003 through November 2004.

{[}2{]} Sex, SLE, and Signal Transduction (National Institutes of
Health, R01 AR50878, Principal Investigator: Jill Jacobson). Total
direct costs of \$250,000 in 2002. The major goal of this project is to
demonstrate that the increase in immune responsiveness to
gonadotropin-releasing hormone in female mice plays a pivotal role in
increased incidence of autoimmune disease. I was a co-investigator
funded at 5\% effort from December 2001 through November 2006.

{[}1{]} Pediatric Pharmacology Research Unit (National Institutes of
Health, U01 HD31313, Principal Investigator: Greg Kearns). Total direct
costs of \$361,932 in 2002. The major goal of this project is to conduct
basic and translational research focusing on regulatory events in
pharmacogenetics, the ontogeny of important drug metabolizing enzymes,
and pharmacodynamics that will advance our understanding of age-related
differences in drug disposition and action. I was a co-investigator
funded at 3\% effort from May 1996 to December 2003.

\subsection{List of awards}\label{list-of-awards}

{[}11{]} Honorable mention in the R programming contest sponsored by
Revolution Analytics, for R functions to predict the completion date for
a prospective clinical trial, May 2014.

{[}10{]} Co-author of a research presentation (Nopper A, Wright T, Tee
R, Horii K, Simon S, Popovic J, Alon U. Bone Density in Children With
Hemangiomas Treated With Systemic Glucocorticoids) which received the
ISSVA ``R Schobinger Award'' for best clinical paper presented during
the 16th International Workshop on Vascular Anomalies, June 2006,
Milano, Italy.

{[}9{]} Co-author of a research paper (Miller-Hansen DR, Nelson PB,
Widen JE, Simon SD. Am J Audiology 2003;24(1);16-18) which received the
Editor's Award for the American Journal of Audiology for the most
outstanding publication in the calendar year 2003.

{[}8{]} Co-inventor on U.S. Patent, A System and Method for Monitoring
and Analyzing Data Trends of Interest Within an Organization, filed
December 2000, and sold to Cerner Corporation, February 2008.

{[}7{]} Co-author of a research paper (Moorman WJ, Skaggs SR, Clark JC,
Turner TW, Sharpnack DD, Murrell JA, Simon SD, Chapin RE, Schrader SM
{[}1998{]}. Male reproductive effects of lead, including species
extrapolation for the rabbit model. Reproductive Toxicology 12(3):
333-46.) which was the winner of the Alice Hamilton Award for 1999 for
the best research paper in occupational safety and health (Biological
Sciences Category).

{[}6{]} My presentation, ``Medical Statistics Case Studies on the Web,''
was voted as the best presentation in the area of Teaching Statistics in
the Health Sciences at the Joint Statistical Meetings in Anaheim CA,
August 1997.

{[}5{]} Received Public Health Service Commendation Medal in July 1994
for my oversight of a laboratory modernization effort that included the
use of a client-server data base system, remote monitoring of laboratory
systems, automation of chemical workups using laboratory robotics,
laptop systems for data collection in the field, and real-time data
collection interfaces for Liquid and Gas Chromatography instruments.

{[}4{]} Co-author of a research paper (Stettler LE, Platek SF, Riley RD,
Mastin JP, Simon SD {[}1991{]}. Lung particulate burdens of subjects
from Cincinnati, Ohio urban area. Scanning Microscopy 5(1):85-94) which
was a finalist for the Alice Hamilton Award for 1992 for the best
research paper in occupational safety and health.

{[}3{]} Co-author of a research paper (Schrader SM, Turner TW, Simon SD
{[}1990{]}. Longitudinal Study of Semen Quality of Unexposed Workers:
II. Sources of Variation of Sperm Head Morphometry. Journal of Andrology
11(1):32-39) which was the winner of the Alice Hamilton Award for 1991
for the best research paper in occupational safety and health. This
paper was nominated by NIOSH for the Charles C. Shepard Science Award
for best scientific research paper in CDC.

{[}2{]} Received Public Health Service Achievement Medal in April 1990
for producing cost savings and greater operational efficiency through
the use of micro-computer analysis of statistical data.

{[}1{]} Co-author of a research paper (Schrader SM, Turner TW,
Breitenstein MJ, Simon SD {[}1988{]}. Longitudinal Study of Semen
Quality of Unexposed Workers. I. Study Overview. Reproductive Toxicology
2:183-190) which was the winner of the Alice Hamilton Award for 1990 for
the best research paper in occupational safety and health.

\subsection{Professional organizations and special committee
assignments}\label{professional-organizations-and-special-committee-assignments}

{[}34{]} Judge, Health Sciences Student Research Summit, March 2023. I
was assigned four posters to review for the Health Sciences Student
Research Summit, a conference to highlight research done by local
students.

{[}33{]} Member, Non-tenure-track, non-physician promotion committee
(April 2023 to present). This group reviews promotion applications to
Associate and Full Professor for non-tenure-track, non-physician faculty
in the School of Medicine. This includes faculty with joint appointments
at various allied hospitals.

{[}32{]} Member, Leadership book club (July 2022 to present). This group
of faculty in the UMKC School of Medicine meets monthly to discuss books
relevant to our work.

{[}31{]} Member, Research oversight, School of Medicine, St.~Joseph
Campus (November 2022 to present). This is an informal working group,
headed by Kathleen Spears, to advise medical students at the St.~Joseph
Campus on various research projects.

{[}30{]} President of the Kansas City R Users Group (January 2014 to
present). This group meets monthly to share information and experiences
with the R programming language. When I took over, the group was
threatening to fold because meetings had as few as three or four people
attending. I converted the group's publicity from a proprietary email
notification system to meetup.com, changed the format of the meetings to
alternate between beginner and advanced sessions, found an alternate
meeting site with better seating and audiovisual equipment, and
recruited new speakers for the meetings. The group has increased its
membership and now has 10 to 20 attendees at most meetings.

{[}29{]} Co-organizer (with M Gerkovich) of Secondary Data Analysis
working group (January 2014 through June 2015). This group meets monthly
and members take turns sharing information about interesting data
sources they have used. In the fall semester of 2015, this group was
folded in with the graduate student seminar.

{[}28{]} Member of the Complementary and Alternative Medicine (CAM)
Statisticians committee within the American Statistical Association
(ASA) (August 2012 to August 2016). This group meets annually at the
Joint Statistical Meetings and conducts additional business by email and
telephone. The members share experiences and concerns about statistical
issues associated with CAM research studies, learns about research
developments at the National Center for Complementary and Integrative
Health, and provides outreach to CAM practitioners who wish to study
their therapies in a rigorous and controlled fashion.

{[}27{]} Provided feedback on presentations at the Greater Kansas City
Society of Health System Pharmacy Resident Research Day, September 2014,
September 2015, and September 2017.

{[}26{]} Member of the ASA Traveling Course committee (August 2004 to
August 2006). The Traveling Course committee solicits speakers who are
willing to offer a one day short course on an important topic in
Statistics. This course is available for three local ASA chapters. The
committee reviews applications and selects three chapters to host the
course. The committee then arranges all of the travel and other
logistics along with the appropriate local ASA chapter officers.

{[}25{]} Member of the Kansas University Medical Center Data Safety and
Monitoring Executive Committee (January 2002 to June 2008) and
Vice-Chair (January 2005 to June 2008). This committeeprovides on-going
review of research projects at the Kansas University Medical Center
where there is a conflict of interest for the primary investigator or
where early stopping of a trial would be considered for early evidence
of efficacy or early evidence of an unacceptable rate of adverse events.

{[}24{]} Member of the University of Missouri Kansas City School of
Medicine search committee for a M.S. Biostatistics position (April 2007
to June 2007). Assisted with review of resumes, personal interviews, and
evaluation.

{[}23{]} Member of the University of Missouri Kansas City School of
Medicine search committee for a Ph.D.~Biostatistics position (July 2006
to April 2007). Assisted with review of resumes, personal interviews,
and evaluation.

{[}22{]} Member of the University of Missouri Kansas City School of
Nursing search committee for a Ph.D.~Biostatistics position (September
2005 to November 2005). Assisted with personal interviews and
evaluation.

{[}21{]} Chaired the planning committee for the April 2004 short course
by C Davis on Statistical Methods for the Analysis of Repeated
Measurements. Served on the planning committees for: the February 1996
short course by R. Helms on Mixed Models, the March 1993 short course by
R. Hogg on Total Quality Management; the April 1992 U.S. Public Health
Service Professional meeting; the March 1992 ENAR Biometrics meeting;
and the November 1990 Ohio Statistics Conference.

{[}20{]} Council of Chapters Representative (July 2001 to June 2004),
President (July 1998 to June 2000), and Vice President (July 1997 to
June 1998) for the Kansas/Western Missouri chapter of the ASA. The
Kansas-Western Missouri chapter provides statistics seminars and other
presentations from local and national experts for locations in Western
Missouri and in Kansas. As Council of Chapters representative, I served
as a liaison between the local ASA chapter and the national office. As
president, I provided support and guidance to all the other officers of
the chapter and represented the club to the parent organization (the
American Statistical Association). As vice-president, I recruited
speakers and organized the time and location of statistics
presentations.

{[}19{]} Served on an advisory committee to develop a Health Outcomes
Research Masters Degree program for the University of Missouri, Kansas
City, October 1999 to February 2000.

{[}18{]} Founder and president of the Childrens Mercy Hospital Internet
Users Group, April 1997 to June 1998. The Internet Users Group provides
a forum for individuals to share their expertise on Internet topics. The
steering committee identifies speakers and topics for meetings and
prepares publicity. As founder, I scheduled and ran an organizational
meeting for the group, developed the structure and goals of this group,
and recruited volunteers for the steering committee. As a president of
the steering committee, I arrange meetings of this group and provide
overall guidance.

{[}17{]} Member of the Division Quality Council, July 1995 to May 1996.
The Quality Council identifies quality problems within a division of 100
employees at NIOSH, recruits teams to address these problems, and
provides continual oversight and feedback to these teams.

{[}16{]} Secretary/Treasurer of Cincinnati Chapter of the American
Statistical Association, July 1995 to June 1996. The Cincinnati chapter
provides statistics seminars and other presentations from local and
national experts for locations in and around Cincinnati, including
Oxford, Ohio. As Secretary/Treasurer, I obtained meeting rooms and
appropriate audiovisual equipment for Statistics presentations, prepared
publicity for each meeting, maintained the organization's mailing list,
and balanced the organization's checking account.

{[}15{]} Member of the Quality Management Advisory Board for the College
of Mount St.~Joseph, May 1995 to May 1996. The Quality Management
Advisory Board provides advice to the College of Mount St.~Joseph
regarding its newly developed undergraduate degree program in Quality
Management. The board consists of representatives of the Cincinnati
business and government communities as well as students and faculty from
the Quality Management program itself.

{[}14{]} Co-chair of the Statistics, Data Analysis and Modeling Section
of the 1995 Midwest SAS Users Group conference in Cleveland OH (with D
Wild). The co-chair provides review of all papers submitted to this
section of the conference. I contacted presenters to verify audiovisual
needs and met with presenters prior to the start of the meeting to
answer questions. In addition, I performed any necessary troubleshooting
before or during the start of each session at the conference, and
ensured that evaluation forms were completed and delivered to the
conference organizer at the end of each session.

{[}13{]} Cincinnati NIOSH Representative to the Cincinnati Federal
Executive Board Reinvention Subcommittee, October 1994 to September
1995. The reinvention subcommittee provides Cincinnati Federal
Government agencies with a professional source of expertise on the
issues addressed in the National Performance Review, and presents
various seminars that address handling change and coping with
streamlining initiatives.

{[}12{]} Team Leader for NIOSH Organizational Climate Survey Team (July
1993 to January 1994), and team member (January 1994 to May 1996). This
Organizational Climate Survey Team distributed and analyzed a survey in
August and September 1993 to assess the perceptions of NIOSH employees
towards Total Quality Management. Our team presented a report detailing
our analysis of this survey to the NIOSH Quality Council in December
1993. The team has proposed a revision of this survey, which we plan to
implement on a semi-regular basis (e.g., every two to three years) to
assess the changes in employee perceptions over time.

{[}11{]} Team Facilitator for Division Manuscript Review Team, June 1993
to May 1996. The Manuscript Review Team applied Total Quality Management
principles to implement improvements to the manuscript review procedure
for a division of 100 employees at NIOSH. Our goal was to decrease the
time required for internal review while still maintaining quality of the
manuscripts. As facilitator, I did not participate in the actual
decisions being made, but instead worked to ensure that decisions made
are founded on team consensus and full, open, and honest participation
from all team members. I encouraged the use of TQM tools such as flow
charts and cause-and-effect diagrams.

{[}10{]} Member of NIOSH Surveillance Coordinating Group, January 1993
to May 1996. The Surveillance Coordinating Group coordinates all
surveillance related research throughout NIOSH, prepares reports on
surveillance activities for the public, and sets standards on variable
naming definitions.

{[}9{]} Member of CDC Statistical Advisory Group, October 1992 to May
1996. The Statistical Advisory Group advises the Assistant Director for
Science of the Centers for Disease Control and Prevention on matters of
policy relating to Statistics.

{[}8{]} Chair of the Division Seminar Series committee, September 1992
to August 1994. The Seminar Series committee organizes weekly research
and safety seminars for a division of 100 employees at NIOSH, and
prepares publicity for these seminars.

{[}7{]} Assisted in the preparation of a poster ``Statistics at NIOSH,''
which was on display at the Hamilton County Public Library (main branch)
during November 1989 to commemorate the 150th anniversary of the
American Statistical Association.

{[}6{]} Founder and Co-chair of NIOSH Statistical Discussion Group, July
1989 to May 1996. The Statistics Discussion Group provides a forum for
informal discussion of professional issues relevant to Statisticians in
the five Cincinnati Divisions of NIOSH and the two Morgantown Divisions.
As founder of the Statistical Discussion Group, I scheduled and ran an
organizational meeting for the group, developed the structure and goals
of this group, and recruited a co-chair for the group. As co-chair, I
was responsible for organizing meetings and for publicity for
Statisticians in two Divisions at the Taft Laboratories in Cincinnati.

{[}5{]} Organizer and chair of the executive board for the Division
Microcomputer Users Group, September 1988 to March 1989, and member of
the executive board, April 1989 to June 1990. The Microcomputer Users
Group coordinated microcomputer purchases and provided educational
services for a division of 100 employees at NIOSH. As founder of this
group, I scheduled and ran an organizational meeting, developed the
structure and goals of this group, and recruited volunteers for the
steering committee. I was responsible for coordinating publicity and
monthly meetings. I served as the liaison to the Hamilton Users Group, a
similar group for two other divisions in NIOSH.

{[}4{]} Charter President of the Bowling Green IBM PC Users Group,
September 1985 to May 1987. I was responsible for recruiting speakers
for their bi-monthly meetings, for publicizing these meetings, and for
maintaining their membership list. I supervised the public domain
librarian and the SYSOPs of the group's Bulletin Board System.

{[}3{]} President (June 1985 to May 1986); Vice President (June 1984 to
May 1985); and Secretary/Treasurer (June 1983 to May 1984) of the
Northwest Ohio chapter of the American Statistical Association. The
Northwest Ohio chapter provides statistics seminars and other
presentations from local and national experts for locations in Northwest
Ohio. As president, I provided support and guidance to all the other
officers of the club and represented the club to the parent organization
(the American Statistical Association). As vice-president, I recruited
speakers and organized the time and location of statistics
presentations. As secretary/treasurer, I obtained meeting rooms and
appropriate equipment for statistical presentations, prepared publicity
for each meeting, maintained the organization's mailing list, and
balanced the checking account.

{[}2{]} Programmed and tested in 1987 with LM Quinn a set of Turbo
Pascal programs (SAMPLE, SAMSTAT, CONFINT, CHART, and HISTGRAM) that
illustrate the random sampling process on real world data sets. In my
Business Statistic courses, I have run these programs as part of my
lectures on the Central Limit Theorem, on the effect of sample size on
sampling error, and on the meaning of confidence in a 95\% confidence
interval. In addition, I have used these programs in a presentation,
``Simulating Random Sampling'' (see above) and in ``Using some Turbo
Pascal programs to simulate random sampling in the classroom'' (see
above).

{[}1{]} Programmed, documented, and tested in 1986 a series of
benchmarks written in BASIC and Pascal (NUMACC.BAS and NUMACC.PAS) to
evaluate numerical accuracy.

\subsection{List of computer skills}\label{list-of-computer-skills}

\begin{itemize}
\tightlist
\item
  Bibliographic software: EndNotes, Knowledge Finder, Mendeley, Zotero.
\item
  Cloud storage: Box, Dropbox, iCloud, OneDrive.
\item
  Database software: Microsoft Access, Microsoft SQL Server, Oracle,
  PC-File, SQLite.
\item
  Electronic Health Records software: i2b2.
\item
  Graphics software: ACDSee, Metafile Companion, Photoshop Elements,
  SigmaPlot.
\item
  Internet systems: File Transfer Protocol, Gopher, Telnet, USENET,
  WordPress, World Wide Web.
\item
  Mathematical software: MathCAD, Mathematica, MathType.
\item
  Operating systems: Linux (Raspbian), MS-DOS, OS/2, Windows
\item
  Presentation software: Powerpoint.
\item
  Programming languages: BASIC, C, C++, FORTRAN, Pascal, Perl, PL/1,
  Python, Visual BASIC.
\item
  Spreadsheets: Excel, Lotus 1-2-3, SuperCalc.
\item
  Statistical software: AMOS, BMDP, IMSL, JMP, LogXact, MINITAB, nQuery
  Advisor, OpenBUGS, R (including RStudio and tidyverse), RATS, S-Plus,
  S-Plus/Wavelets, SAS (including SAS University), Stan, SPIDA, SPSS,
  STATA, Statgraphics, StatXact, Systat, WinBUGS.
\item
  Utility software: DBMS/COPY, Norton Anti-virus, Notepad++, TextPad,
  WinZip.
\item
  Word processing: LaTeX (MIKTeX), RMarkdown (including blogdown,
  bookdown, and pagedown), Word, Word Perfect.
\item
  Other experience: I have supervised the implementation of a Novell
  Local Area Network and Microsoft SQL Server, a client-server database,
  both for a division of 100+ employees. I have authored web pages using
  Adobe Acrobat, Adobe Dreamweaver, Camtasia Studio, and Microsoft Front
  Page. I have used and taught IBM mainframe Job Control Language,
  including tape and disk I/O.
\end{itemize}

\subsection{List of volunteer efforts}\label{list-of-volunteer-efforts}

{[}11{]} Head of the Second Friday Movie Group, December 2004 to
present, a group that offers a social outlet formerly just for members
of St.~Pauls United Methodist Church in Lenexa, but now open to anyone.
We meet on the second Friday of each month to see a first run movie and
afterwards visit a nearby restaurant to discuss the movie. As head of
the group, I maintain a mailing list of members, select an appropriate
movie, announce the choice, and collect RSVPs. The Movie Group hosts an
Academy Awards party every year.

{[}10{]} Performer with the Back Porch Cloggers, December 1999 to May
2006, a group that performed Appalachian Clogging. We would regularly
entertain at various local retirement homes and nursing facilities.

{[}9{]} Served as Vice President Education (July 2001 to June 2003),
President (July 2000 to June 2001) and Sergeant-at-Arms (February 2000
to June 2000) for the Bluejacket Toastmasters Club. Toastmasters
provides a mutually supportive and positive learning environment for the
development of communication and leadership skills. As Vice President
Education, I prepared a program of speakers and evaluators for each
meeting and reviewed credentials of members applying for Competent
Toastmaster, Able Toastmaster, and Distinguished Toastmaster
certifications. As president, I provided support and guidance to all the
other officers of the club and represent the club to the parent
organization (Toastmasters International) and to outside groups. As
Sergeant-at-Arms, I prepared the room for each meeting and provided
refreshments. I was selected as the winner of the speech contest for my
speech ``I married a dog nut,'' March 2003, humorous speech contest for
my speech ``The one thing that can gross out a doctor,'' September 2001,
and also won at the Area K-1 contest, and the Division K contest.

{[}8{]} Tutored for the Adult Basic Literacy program, March 2000 to
December 2001, two hours per week. I provided assistance to adults who
were studying for their GED degree.

{[}7{]} Tutored students at Longfellow Elementary School, November 1999
to March 2000, one hour per week. I provided assistance to a third grade
student in various academic areas, mostly math.

{[}6{]} Edited the Red Oak Hills Homeowners Association newsletter,
November 1998 to April 2002. This newsletter appears four or more times
a year and highlights information of interest to the 275 home owners in
Red Oak Hills.

{[}5{]} Served as President (July 1994 to June 1995) and charter member
(March 1993 to May 1996) of Cincinnati NIOSH Toastmasters. Toastmasters
provides a mutually supportive and positive learning environment for the
development of communication and leadership skills. As president, I
provided support and guidance to all the other officers of the club and
represented the club to the parent organization (Toastmasters
International) and to outside groups. In April 1995, I earned the
Competent Toastmasters Certificate for successful completion of the ten
speech projects in the Toastmasters Communication and Leadership
Program. I was selected as the winner of the humorous speech contest for
my speech ``Puns and Poetry,'' September 1994, and also won at the Area
32 contest. I was selected as winner of the Table Topics (extemporaneous
speech) contest, September 1994 for my response to the question ``What
famous person would you like to have lunch with?''

{[}4{]} Keyperson for the Combined Federal Campaign, Fall 1989. I passed
out pledge forms for the Combined Federal Campaign, a program for
federal employees similar to United Way, and answered questions about
this program.

{[}3{]} Assisted in the proctoring and grading of tests in a Math
contest sponsored by Miami University, Middletown, March 1988.

{[}2{]} Selected as faculty advisor for Bowling Green State University
Chess Club, Fall 1984 to Spring 1987.

{[}1{]} Computerized the mailing list with DK Wild for the Wood County
Humane Society using PC-FILE III, Spring 1985.

\subsection{Hobbies and Interests}\label{hobbies-and-interests}

\begin{itemize}
\tightlist
\item
  Appalachian Clogging, Bridge, Cats, Chess, Computers, Contra Dancing,
  Crossword Puzzles, Languages, Running (5K, 4 mile, and 10K races),
  Travel.
\end{itemize}

\subsection{List of publications}\label{list-of-publications}

The definitive list of publications is available at
\url{https://www.ncbi.nlm.nih.gov/myncbi/browse/collection/48071026/?sort=date&direction=descending}

{[}122{]} Berkley-Patton J et al.~HIV Testing in African American
Churches: Results from the Taking It to the Pews Cluster-Randomized
Trial, submitted to the American Journal of Public Health.

{[}121{]} Wilson, Neal J Allenbrand R, Friedman E, Kennedy K, Roberts A,
Simon S. Visualizing Parcel-Level Lead Risk Using an Exterior
Housing-Based Index. International Journal of Environmental Research and
Public Health. 2024 Dec 27, 22(1), 16. \url{doi:10.3390/ijerph22010016}

{[}120{]} Pradeep Kumar D, Zanotto T, Cozart JS, Bruce AS, Befort C,
Siengsukon C, Shook R, Lynch S, Mahmoud R, Simon S, Hibbing PR, Drees B,
Huebner J, Bradish T, Robichaud J, Sosnoff JJ, Bruce JM. Association
between frailty and sleep quality in people living with multiple
sclerosis and obesity: An observational cross-sectional study. Mult
Scler Relat Disord. 2024 Jan;81:105154. doi:
10.1016/j.msard.2023.105154.

{[}119{]} Bruce JM, Cozart JS, Shook RP, Befort C, Siengsukon CF, Simon
S, Lynch SG, Mahmoud R, Drees B, Posson P, Hibbing PR, Huebner J,
Bradish T, Robichaud J, Bruce AS. Modifying diet and exercise in
multiple sclerosis (MoDEMS): A randomized controlled trial for
behavioral weight loss in adults with multiple sclerosis and obesity.
Mult Scler. 2023 Dec;29(14):1860-1871. doi: 10.1177/13524585231213241.

{[}118{]} Cozart JS, Bruce AS, Shook RP, Befort C, Siengsukon C, Simon
S, Lynch SG, Mahmoud R, Drees B, Posson P, Hibbing PR, Huebner J,
Bradish T, Robichaud J, Bruce JM. Body metrics are associated with
clinical, free-living, and self-report measures of mobility in a cohort
of adults with obesity and multiple sclerosis. Mult Scler Relat Disord.
2023 Nov;79:105010. doi: 10.1016/j.msard.2023.105010. Epub 2023 Sep 15.

{[}117{]} Cozart JS, Bruce AS, Befort C, Siengsukon C, Lynch SG, Punt S,
Simon S, Shook RP, Huebner J, Bradish T, Robichaud J, Bruce JM. A pilot
study evaluating the prefeasibility of a behavioral weight loss program
in people with multiple sclerosis. Prev Med Rep.~2023 Sep 22;36:102437.
doi: 10.1016/j.pmedr.2023.102437.

{[}116{]} Derose KP, Berkley-Patton J, Hamilton-Burgess C, Thompson CB,
Williams ED, Simon S, Allsworth JE. Correlates of HIV-Related Stigmas
Among African American Church-Affiliated Populations in Kansas City.
AIDS Educ Prev. 2023 Feb;35(1):54-68. doi: 10.1521/aeap.2023.35.1.54.

{[}115{]} Wilson NJ, Friedman E, Kennedy K, Manolakos PT, Reierson L,
Roberts A, Simon S. Using exterior housing conditions to predict
elevated pediatric blood lead levels. Environ Res. 2023 Feb
1;218:114944. doi: 10.1016/j.envres.2022.114944.

{[}114{]} Berkley-Patton J, Thompson CB, Templeton T, Burgin T, Derose
KP, Williams E, Thompson F, Catley D, Simon SD, Allsworth JE. COVID-19
Testing in African American Churches Using a Faith-Health- Academic
Partnership. Am J Public Health. 2022 Nov;112(S9):S887-S891. doi:
10.2105/AJPH.2022.306981.

{[}113{]} Wamkpah N, Shrestha A, Szalzman G, Simon S, Suman S, Poisner
A, Molteni A. Renin-Angiotensin Blockade Reduces Readmission for Acute
Chest Syndrome in Sickle Cell Disease. Cureus 2022-03-28, 14(3): e23567,
doi: 10.7759/cureus.23567.

{[}112{]} Bruce JM, Cozart JS, Shook RP, Ruppen S, Siengsukon C, Simon
S, Befort C, Lynch S, Mahmoud R, Drees B, Norouzinia AN, Bradish T,
Posson P, Hibbing PR, Bruce AS (2021). Modifying Diet and Exercise in MS
(MoDEMS): Study design and protocol for a telehealth weight loss
intervention for adults with obesity \& Multiple Sclerosis. Contemporary
Clinical Trials. 2021-08-01; 107: 106495.

{[}111{]} Dubin JR, Simon SD, Norrell K, Perera J, Gowen J, Cil A. Risk
of Recall Among Medical Devices Undergoing US Food and Drug
Administration 510(k) Clearance and Premarket Approval, 2008-2017. JAMA
Network Open. 2021-05-06; 4(5): e217274.
\url{doi:10.1001/jamanetworkopen.2021.7274}

{[}110{]} Simon S. A Peek into the Inner Workings of Pandemic Prediction
Models. Mo Med. 2021 May-Jun;118(3):259-263.

{[}109{]} Younis M, Elkaryoni A, Williams GW, Jakhar I, Suman S, Simon
S, Salzman G. The Use of Direct Oral Anticoagulants in the Management of
Venous Thromboembolism in Patients With Obesity. Cureus 2020-08-25,
12(8): e10006. \url{doi:10.7759/cureus.10006.50/35}

{[}108{]} Simon SD. Centers for Disease Control and Prevention. In
Encyclopedia of Big Data. Schintler LA, McNeely CL (eds.). In press.

{[}107{]} Simon SD. Census Bureau (U.S.). In Encyclopedia of Big Data.
Schintler LA, McNeely CL (eds.). In press.

{[}106{]} Berkley-Patton J, Bowe Thompson C, Goggin K, Catley D, Berman
M, Bradley-Ewing A, Derose K, Resnicow K, Allsworth J, \& Simon S. A
religiously-tailored, multilevel intervention in African American
churches to increase HIV testing rates: rationale and design of the
Taking It to the Pews cluster randomized trial. Contemporary Clinical
Trials 2019; 86: 105848.

{[}105{]} Simon SD. Getting Paid: How to Determine Your Fee. Amstat
News. 2017 (August 1). This article was not peer-reviewed.

{[}104{]} Simon SD. Why Be an Independent Consultant? Amstat News. 2017
(April 1). This article was not peer-reviewed.

{[}103{]} Barker JP, Simon SD, Dubin J. The Methodology of Clinical
Studies Used by the FDA for Approval of High-Risk Orthopaedic Devices.
The Journal of Bone and Joint Surgery. American Volume. 2017;
99(9):711-719. PMID: 28463914

{[}102{]} Parthiban A, Levine JC, Nathan M, Marshall JA, Shirali GS,
Simon SD, Colan SD, Newburger JW, Raghuveer G. Implementation of a
Quality Improvement Bundle Improves Echocardiographic Imaging after
Congenital Heart Surgery in Children. Journal of the American Society of
Echocardiography : official publication of the American Society of
Echocardiography. 2016; 29(12):1163-1170.e3. PMID: 27742240

{[}101{]} Jiang Y, Guarino P, Ma S, Simon S, Mayo MS, Raghavan R,
Gajewski BJ. Bayesian accrual prediction for interim review of clinical
studies: open source R package and smartphone application. Trials. 2016;
17(1):336. PMID: 27449769, PMCID: PMC4957321

{[}100{]} Prindaville B, Simon SD, Horii KA. Dermatology-related
outpatient visits by children: Implications for workforce and pediatric
education. Journal of the American Academy of Dermatology. 2016;
75(1):228-9. PMID: 27317526

{[}99{]} Berkley-Patton J, Thompson CB, Moore E, Hawes S, Simon S,
Goggin K, Martinez D, Berman M, Booker A. An HIV Testing Intervention in
African American Churches: Pilot Study Findings. Annals of Behavioral
Medicine : a publication of the Society of Behavioral Medicine. 2016;
50(3):480-5. NIHMSID: NIHMS752446 PMID: 26821712, PMCID: PMC5026504

{[}98{]} Baker JM, Pace HA, Ladesich JB, Simon SD. Evaluation of the
Impact of Corticosteroid Dose on the Incidence of Hyperglycemia in
Hospitalized Patients with an Acute Exacerbation of Chronic Obstructive
Pulmonary Disease. Hospital Pharmacy. 2016; 51(4):296-304. PMID:
27303077, PMCID: PMC4896332

{[}97{]} Parthiban A, Levine JC, Nathan M, Marshall JA, Shirali GS,
Simon SD, Colan SD, Newburger JW, Raghuveer G. Impact of Variability in
Echocardiographic Interpretation on Assessment of Adequacy of Repair
Following Congenital Heart Surgery: A Pilot Study. Pediatric Cardiology.
2016; 37(1):144-50. PMID: 26358473

{[}96{]} Anderst J, Teran P, Dowd MD, Simon S, Schnitzer P. The
association of the rapid assessment of supervision scale score and
unintentional childhood injury. Child Maltreatment. 2015; 20(2):141-5.
PMID: 25601937

{[}95{]} Knapp JF, Simon SD, Sharma V. Does active dissemination of
evidence result in faster knowledge transfer than passive diffusion?: An
analysis of trends of the management of pediatric asthma and croup in US
emergency departments from 1995 to 2009. Pediatric Emergency Care. 2015;
31(3):190-6. PMID: 24694945

{[}94{]} Jiang Y, Simon S, Mayo MS, Gajewski BJ. Modeling and validating
Bayesian accrual models on clinical data and simulations using adaptive
priors. Statistics in Medicine. 2015; 34(4):613-29. NIHMSID: NIHMS637332
PMID: 25376910, PMCID: PMC4314351

{[}93{]} Simon SD. R for Big Data Analysis, a chapter in Big Data
Analysis for Bioinformatics and Biomedical Discoveries. Ye SQ, editor.
New York, NY: Chapman \& Hall; 2015.

{[}92{]} Patterson ME, Marken P, Zhong Y, Simon SD, Ketcherside W.
Comprehensive electronic medical record implementation levels not
associated with 30-day all-cause readmissions within Medicare
beneficiaries with heart failure. Applied Clinical Informatics. 2014;
5(3):670-84. PMID: 25298808, PMCID: PMC4187085

{[}91{]} Berkley-Patton JY, Moore E, Berman M, Simon SD, Thompson CB,
Schleicher T, Hawes SM. Assessment of HIV-related stigma in a US
faith-based HIV education and testing intervention. Journal of the
International AIDS Society. 2013; 16(3 Suppl 2):18644. PMID: 24242259,
PMCID: PMC3833192

{[}90{]} Shah A, Kumar S, Simon SD, Singh DP, Kumar A. HIV gp120- and
methamphetamine-mediated oxidative stress induces astrocyte apoptosis
via cytochrome P450 2E1. Cell Death \& Disease. 2013; 4:e850. PMID:
24113184, PMCID: PMC3824683

{[}89{]} Knapp JF, Simon SD, Sharma V. Variation and trends in ED use of
radiographs for asthma, bronchiolitis, and croup in children.
Pediatrics. 2013; 132(2):245-52. PMID: 23878045

{[}88{]} Gajewski BJ, Simon SD, Carlson SE. On the Existence of Constant
Accrual Rates in Clinical Trials and Direction for Future Research.
International Journal of Statistics and Probability. 2012; 1(2):p43.
NIHMSID: NIHMS389648 PMID: 23869201, PMCID: PMC3712523

{[}87{]} Patterson ME, Marken PA, Simon SD, Hackman JL, Schaefer RS.
Associations between the concurrent use of clinical decision support and
computerized provider order entry and the rates of appropriate
prescribing at discharge. Applied Clinical Informatics. 2012;
3(2):186-96. PMID: 23646071, PMCID: PMC3613019

{[}86{]} Simon SD. Tackling conflicts of interest. Subjective editorials
and clinical reviews require proof of objectivity. BMJ (Clinical
Research Ed.). 2011; 343:d5601. PMID: 21896589

{[}85{]} Gilman AL, Jacobsen C, Bunin N, Levine J, Goldman F, Bendel A,
Joyce M, Anderson P, Rozans M, Wall DA, Macdonald TJ, Simon S, Kadota
RP. Phase I study of tandem high-dose chemotherapy with autologous
peripheral blood stem cell rescue for children with recurrent brain
tumors: a Pediatric Blood and MarrowTransplant Consortium study.
Pediatric Blood \& Cancer. 2011; 57(3):506-13. PMID: 21744474

{[}84{]} Talib N, Onikul R, Filardi D, Simon S, Sharma V. Effective
educational instruction in preventive oral health: hands-on training
versus web-based training. Pediatrics. 2010; 125(3):547-53. PMID:
20123769

{[}83{]} Mercer AM, Teasley SL, Hopkinson J, McPherson DM, Simon SD,
Hall RT. Evaluation of a breastfeeding assessment score in a diverse
population. Journal of Human Lactation : official journal of
International Lactation Consultant Association. 2010; 26(1):42-8. PMID:
19759350

{[}82{]} Miller MK, Denise Dowd M, Gratton MC, Cai J, Simon SD.
Pediatric out-of-hospital emergency medical services utilization in
Kansas City, Missouri. Academic Emergency Medicine : official journal of
the Society for Academic Emergency Medicine. 2009; 16(6):526-31. PMID:
19426299

{[}81{]} Midyett LK, Grunt J, Simon SD. Noninvasive radial artery
tonometry augmentation index and urinary albumin/creatinine levels in
early adolescents with type 1 diabetes mellitus. Journal of Pediatric
Endocrinology \& Metabolism : JPEM. 2009; 22(6):531-7. PMID: 19694200

{[}80{]} Jones BL, Abdel-Rahman SM, Simon SD, Kearns GL, Neville KA.
Assessment of histamine pharmacodynamics by microvasculature response of
histamine using histamine iontophoresis laser Doppler flowimetry.
Journal of Clinical Pharmacology. 2009; 49(5):600-5. PMID: 19318525

{[}79{]} Knapp JF, Simon SD, Sharma V. Quality of care for common
pediatric respiratory illnesses in United States emergency departments:
analysis of 2005 National Hospital Ambulatory Medical Care Survey Data.
Pediatrics. 2008; 122(6):1165-70. PMID: 19047229

{[}78{]} Graf WD, Kayyali HR, Alexander JJ, Simon SD, Morriss MC.
Neuroimaging-use trends in nonacute pediatric headache before and after
clinical practice parameters. Pediatrics. 2008; 122(5):e1001-5. PMID:
18838461

{[}77{]} Schurman JV, Danda CE, Friesen CA, Hyman PE, Simon SD, Cocjin
JT. Variations in psychological profile among children with recurrent
abdominal pain. Journal of Clinical Psychology in Medical Settings.
2008; 15(3):241-51. PMID: 19104969

{[}76{]} Gajewski BJ, Simon SD. A One-Hour Training Seminar on Bayesian
Statistics for Nursing Graduate Students. The American Statistician.
2008 August 01; 62(3):190-4.

{[}75{]} Gajewski BJ, Simon SD, Carlson SE. Predicting accrual in
clinical trials with Bayesian posterior predictive distributions.
Statistics in Medicine. 2008; 27(13):2328-40. PMID: 17979152

{[}74{]} Sati L, Ovari L, Bennett D, Simon SD, Demir R, Huszar G. Double
probing of human spermatozoa for persistent histones, surplus cytoplasm,
apoptosis and DNA fragmentation. Reproductive Biomedicine Online. 2008;
16(4):570-9. PMID: 18413066

{[}73{]} Adjei AA, Gaedigk A, Simon SD, Weinshilboum RM, Leeder JS.
Interindividual variability in acetaminophen sulfation by human fetal
liver: implications for pharmacogenetic investigations of drug-induced
birth defects. Birth Defects Research. Part A, Clinical and Molecular
Teratology. 2008; 82(3):155-65. PMID: 18232020

{[}72{]} Gaedigk A, Simon SD, Pearce RE, Bradford LD, Kennedy MJ, Leeder
JS. The CYP2D6 activity score: translating genotype information into a
qualitative measure of phenotype. Clinical Pharmacology and
Therapeutics. 2008; 83(2):234-42. PMID: 17971818

{[}71{]} Soden SE, Garrison CB, Lowry JA, Wasserman GS, Simon SD.
Authors' reply to letter regarding the original article ``24-hour
provoked excretion test for heavy metals in children with autism and
typically developing controls, a pilot study'' Clinical Toxicology. 2008
January; 46(10):1098. (Note: I was not a co-author on the original
article, just on this reply.)

{[}70{]} Horii KA, Simon SD, Liu DY, Sharma V. Atopic dermatitis in
children in the United States, 1997-2004: visit trends, patient and
provider characteristics, and prescribing patterns. Pediatrics. 2007;
120(3):e527-34. PMID: 17766497

{[}69{]} Auron A, Warady BA, Simon S, Blowey DL, Srivastava T, Musharaf
G, Alon US. Use of the multipurpose drainage catheter for the provision
of acute peritoneal dialysis in infants and children. American Journal
of Kidney Diseases : the official journal of the National Kidney
Foundation. 2007; 49(5):650-5. PMID: 17472847

{[}68{]} Auron A, Simon S, Andrews W, Jones L, Johnson S, Musharaf G,
Warady BA. Prevention of peritonitis in children receiving peritoneal
dialysis. Pediatric Nephrology (Berlin, Germany). 2007; 22(4):578-85.
PMID: 17216260

{[}67{]} Hampl SE, Carroll CA, Simon SD, Sharma V. Resource utilization
and expenditures for overweight and obese children. Archives of
Pediatrics \& Adolescent Medicine. 2007; 161(1):11-4. PMID: 17199061

{[}66{]} Abdel-Rahman SM, Simon S, Wright KJ, Ndjountche L, Gaedigk A.
Tracking Trichophyton tonsurans through a large urban child care center:
defining infection prevalence and transmission patterns by molecular
strain typing. Pediatrics. 2006; 118(6):2365-73. PMID: 17142520

{[}65{]} Talebizadeh Z, Simon SD, Butler MG. X chromosome gene
expression in human tissues: male and female comparisons. Genomics.
2006; 88(6):675-681. PMID: 16949791

{[}64{]} Ford SP, Leick-Rude MK, Meinert KA, Anderson B, Sheehan MB,
Haney BM, Leeks SR, Simon SD, Jackson JK. Overcoming barriers to oxygen
saturation targeting. Pediatrics. 2006; 118 Suppl 2:S177-86. PMID:
17079621

{[}63{]} Jackson JK, Ford SP, Meinert KA, Leick-Rude MK, Anderson B,
Sheehan MB, Haney BM, Leeks SR, Simon SD. Standardizing nasal cannula
oxygen administration in the neonatal intensive care unit. Pediatrics.
2006; 118 Suppl 2:S187-96. PMID: 17079622

{[}62{]} Dinakar C, Simon S. The link between the Hippocratic Oath and
evidence-based medicine. Annals of allergy, asthma \& immunology :
official publication of the American College of Allergy, Asthma, \&
Immunology. 2006; 96(4):511-3. PMID: 16680920

{[}61{]} Abanses JC, Dowd MD, Simon SD, Sharma V. Impact of rapid
influenza testing at triage on management of febrile infants and young
children. Pediatric Emergency Care. 2006; 22(3):145-9. PMID: 16628094

{[}60{]} Beasley BW, Simon SD, Wright SM. A time to be promoted. The
Prospective Study of Promotion in Academia (Prospective Study of
Promotion in Academia). eJournal of General Internal Medicine. 2006;
21(2):123-9. PMID: 16336619, PMCID: PMC1484667

{[}59{]} Simon SD. Statistical Evidence in Medical Trials. Oxford UK:
Oxford University Press; 2006.

{[}58{]} Richmond W, Colgan G, Simon S, Stuart-Hilgenfeld M, Wilson N,
Alon US. Random urine calcium/osmolality in the assessment of calciuria
in children with decreased muscle mass. Clinical Nephrology. 2005;
64(4):264-70. PMID: 16240897

{[}57{]} Thompson M, Simon SD, Sharma V, Alon US. Timing of follow-up
voiding cystourethrogram in children with primary vesicoureteral reflux:
development and application of a clinical algorithm. Pediatrics. 2005;
115(2):426-34. PMID: 15687452

{[}56{]} Hilgenfeld MS, Simon S, Blowey D, Richmond W, Alon US. Lack of
seasonal variations in urinary calcium/creatinine ratio in school-age
children. Pediatric Nephrology (Berlin, Germany). 2004; 19(10):1153-5.
PMID: 15309603

{[}55{]} Jackson JK, Biondo DJ, Jones JM, Moor PJ, Simon SD, Hall RT,
Kilbride HW. Can an alternative umbilical arterial catheter solution and
flush regimen decrease iatrogenic hemolysis while enhancing nutrition? A
double-blind, randomized, clinical trial comparing an isotonic amino
acid with a hypotonic salt infusion. Pediatrics. 2004; 114(2):377-83.
PMID: 15286220

{[}54{]} George ME, Sharma V, Jacobson J, Simon S, Nopper AJ. Adverse
effects of systemic glucocorticosteroid therapy in infants with
hemangiomas. Archives of Dermatology. 2004; 140(8):963-9. PMID: 15313812

{[}53{]} Hellerstein S, Berenbom M, DiMaggio S, Erwin P, Simon SD,
Wilson N. Comparison of two formulae for estimation of glomerular
filtration rate in children. Pediatric Nephrology (Berlin, Germany).
2004; 19(7):780-4. PMID: 15071771

{[}52{]} Ekekezie II, Thibeault DW, Simon SD, Norberg M, Merrill JD,
Ballard RA, Ballard PL, Truog WE. Low levels of tissue inhibitors of
metalloproteinases with a high matrix metalloproteinase-9/tissue
inhibitor of metalloproteinase-1 ratio are present in tracheal aspirate
fluids of infants who develop chronic lung disease. Pediatrics. 2004;
113(6):1709-14. PMID: 15173495

{[}51{]} Grunt JA, Midyett LK, Simon SD, Lowe L. When should cranial
magnetic resonance imaging be used in girls with early sexual
development? Journal of Pediatric Endocrinology \& Metabolism : JPEM.
2004; 17(5):775-80. PMID: 15237713

{[}50{]} Murphy WJ, Franks JR, Berger EH, Behar A, Casali JG,
Dixon-Ernst C, Krieg EF, Mozo BT, Royster JD, Royster LH, Simon SD,
Stephenson C. Development of a new standard laboratory protocol for
estimation of the field attenuation of hearing protection devices:
sample size necessary to provide acceptable reproducibility. The Journal
of the Acoustical Society of America. 2004; 115(1):311-23. PMID:
14759024

{[}49{]} Miller-Hansen DR, Nelson PB, Widen JE, Simon SD. Evaluating the
benefit of speech recoding hearing aids in children. American Journal of
Audiology. 2003; 12(2):106-13. PMID: 14964326

{[}48{]} Portnoy JM, Simon SD. Is 3-mm less drowsiness important? Annals
of Allergy, Asthma \& Immunology : official publication of the American
College of Allergy, Asthma, \& Immunology. 2003; 91(4):324-5. PMID:
14582809

{[}47{]} Murphy JV, Torkelson R, Dowler I, Simon S, Hudson S. Vagal
nerve stimulation in refractory epilepsy: the first 100 patients
receiving vagal nerve stimulation at a pediatric epilepsy center.
Archives of Pediatrics \& Adolescent Medicine. 2003; 157(6):560-4. PMID:
12796236

{[}46{]} Santos SR, Carroll CA, Cox KS, Teasley SL, Simon SD, Bainbridge
L, Cunningham M, Ott L. Baby boomer nurses bearing the burden of care: A
four-site study of stress, strain, and coping for inpatient registered
nurses. The Journal of Nursing Administration. 2003; 33(4):243-50. PMID:
12690256

{[}45{]} Sharma V, Dowd MD, Swanson DS, Slaughter AJ, Simon SD.
Influence of the news media on diagnostic testing in the emergency
department. Archives of Pediatrics \& Adolescent Medicine. 2003;
157(3):257-60. PMID: 12622675

{[}44{]} De Jonge CJ, Centola GM, Reed ML, Shabanowitz RB, Simon SD,
Quinn P. Human sperm survival assay as a bioassay for the assisted
reproductive technologies laboratory. Journal of andrology. 2003;
24(1):16-8. PMID: 12514073

{[}43{]} Hall RT, Mercer AM, Teasley SL, McPherson DM, Simon SD, Santos
SR, Meyers BM, Hipsh NE. A breast-feeding assessment score to evaluate
the risk for cessation of breast-feeding by 7 to 10 days of age. The
Journal of Pediatrics. 2002; 141(5):659-64. PMID: 12410194

{[}42{]} Carroll CA, Cox KS, Santos SR, Simon SD. Using standard
desk-top tools to monitor medical error rates. Seminars for Nurse
Managers. 2002; 10(2):95-9. PMID: 12092273

{[}41{]} Grunt JA, Schwartz ID, Simon SD, Howard CP. Growth Hormone
Responsiveness in Two Groups of Children With Significantly Short
Stature. The Endocrinologist. 2002; 12:58-65.

{[}40{]} Sharma V, Dowd MD, Slaughter AJ, Simon SD. Effect of rapid
diagnosis of influenza virus type a on the emergency department
management of febrile infants and toddlers. Archives of Pediatrics \&
Adolescent Medicine. 2002; 156(1):41-3. PMID: 11772189

{[}39{]} Simon SD. Is the randomized clinical trial the gold standard of
research? Journal of Andrology. 2001; 22(6):938-43. PMID: 11700857

{[}38{]} Hellerstein S, Simon SD, Berenbom M, Erwin P, Nickell E.
Creatinine excretion rates for renal clearance studies. Pediatric
Nephrology (Berlin, Germany). 2001; 16(8):637-43. PMID: 11519893

{[}37{]} Simon SD. Understanding the odds ratio and the relative risk.
Journal of Andrology. 2001; 22(4):533-6. PMID: 11451349

{[}36{]} Barnes C, Tuck J, Simon S, Pacheco F, Hu F, Portnoy J.
Allergenic materials in the house dust of allergy clinic patients.
Annals of Allergy, Asthma \& Immunology: official publication of the
American College of Allergy, Asthma, \& Immunology. 2001; 86(5):517-23.
PMID: 11379802

{[}35{]} Simon SD. Interpreting positive studies. Journal of Andrology.
2001; 22(3):358-9. PMID: 11330635

{[}34{]} So NP, Osorio AV, Simon SD, Alon US. Normal urinary
calcium/creatinine ratios in African-American and Caucasian children.
Pediatric Nephrology (Berlin, Germany). 2001; 16(2):133-9. PMID:
11261680

{[}33{]} Simon SD. Interpreting negative studies. Journal of Andrology.
2001; 22(1):13-6. PMID: 11191078

{[}32{]} Sharma V, Simon SD, Bakewell JM, Ellerbeck EF, Fox MH, Wallace
DD. Factors influencing infant visits to emergency departments.
Pediatrics. 2000; 106(5):1031-9. PMID: 11061772

{[}31{]} Grajewski B, Cox C, Schrader SM, Murray WE, Edwards RM, Turner
TW, Smith JM, Shekar SS, Evenson DP, Simon SD, Conover DL. Semen quality
and hormone levels among radiofrequency heater operators. Journal of
Occupational and Environmental Medicine. 2000; 42(10):993-1005. PMID:
11039163

{[}30{]} Hall RT, Simon S, Smith MT. Readmission of breastfed infants in
the first 2 weeks of life. Journal of Perinatology : official journal of
the California Perinatal Association. 2000; 20(7):432-7. PMID: 11076327

{[}29{]} Portnoy J, Landuyt J, Pacheco F, Flappan S, Simon S, Barnes C.
Comparison of the Burkard and Allergenco MK-3 volumetric collectors.
Annals of Allergy, Asthma \& Immunology : official publication of the
American College of Allergy, Asthma, \& Immunology. 2000; 84(1):19-24.
PMID: 10674560

{[}28{]} Gaedigk A, Gotschall RR, Forbes NS, Simon SD, Kearns GL, Leeder
JS. Optimization of cytochrome P4502D6 (CYP2D6) phenotype assignment
using a genotyping algorithm based on allele frequency data.
Pharmacogenetics. 1999; 9(6):669-82. PMID: 10634130

{[}27{]} Kliethermes PA, Cross ML, Lanese MG, Johnson KM, Simon SD.
Transitioning preterm infants with nasogastric tube supplementation:
increased likelihood of breastfeeding. Journal of Obstetric,
Gynecologic, and Neonatal Nursing : JOGNN. 1999; 28(3):264-73. PMID:
10363538

{[}26{]} Srivastava T, Simon SD, Alon US. High incidence of focal
segmental glomerulosclerosis in nephrotic syndrome of childhood.
Pediatric Nephrology (Berlin, Germany). 1999; 13(1):13-8. PMID: 10100283

{[}25{]} Schrader SM, Langford RE, Turner TW, Breitenstein MJ, Clark JC,
Jenkins BL, Lundy DO, Simon SD, Weyandt TB. Reproductive function in
relation to duty assignments among military personnel. Reproductive
Toxicology (Elmsford, N.Y.). 1998; 12(4):465-8. PMID: 9717697

{[}24{]} Moorman WJ, Skaggs SR, Clark JC, Turner TW, Sharpnack DD,
Murrell JA, Simon SD, Chapin RE, Schrader SM. Male reproductive effects
of lead, including species extrapolation for the rabbit model.
Reproductive Toxicology (Elmsford, N.Y.). 1998; 12(3):333-46. PMID:
9628556

{[}23{]} Sivak A, Niemeier R, Lynch D, Beltis K, Simon S, Salomon R,
Latta R, Belinky B, Menzies K, Lunsford A, Cooper C, Ross A, Bruner R.
Skin carcinogenicity of condensed asphalt roofing fumes and their
fractions following dermal application to mice. Cancer Letters. 1997;
117(1):113-23. PMID: 9233840

{[}22{]} Weyandt TB, Schrader SM, Turner TW, Simon SD. Semen analysis of
military personnel associated with military duty assignments.
Reproductive Toxicology (Elmsford, N.Y.). 1996; 10(6):521-8. PMID:
8946566

{[}21{]} Combining reproductive studies of men exposed to
2-ethoxyethanol to increase statistical power Schrader SM, Turner TW,
Ratcliffe JW, Welch LS, Simon SD. Occupational Hygiene. 1996; 2:411-5.

{[}20{]} Schrader SM, Ayers-Cumbo M, Turner TW, Simon SD, Ratcliffe J,
Welch L, Grajewski B, Weyandt T. Semen quality across populations.
Molecular Andrology. 1995; 7:105-21.

{[}19{]} Brown KK, Teass AW, Simon SD, Ward EM. A biological monitoring
method for o-Toluidine and Aniline in urine using high performance
liquid chromatography with electrochemical detection. Applied
Occupational and Environmental Hygiene. 1995; 10(6):557-65.

{[}18{]} Schrader SM, Turner TW, Breitenstein MJ, Simon SD. Measuring
male reproductive hormones for occupational field studies. Journal of
Occupational Medicine : official publication of the Industrial Medical
Association. 1993; 35(6):574-6. PMID: 8331437

{[}17{]} Knecht EA, Moorman WJ, Clark JC, Hull RD, Biagini RE, Lynch DW,
Boyle TJ, Simon SD. Pulmonary reactivity to vanadium pentoxide following
subchronic inhalation exposure in a non-human primate animal model.
Journal of Applied Toxicology : JAT. 1992; 12(6):427-34. PMID: 1452976

{[}16{]} Platek SF, Riley RD, Simon SD. The classification of asbestos
fibres by scanning electron microscopy and computer-digitizing tablet.
The Annals of Occupational Hygiene. 1992; 36(2):155-71. PMID: 1326908

{[}15{]} Stettler LE, Platek SF, Riley RD, Mastin JP, Simon SD. Lung
particulate burdens of subjects from the Cincinnati, Ohio urban area.
Scanning Microscopy. 1991; 5(1):85-92; discussion 92-4. PMID: 1647057

{[}14{]} Schrader SM, Turner TW, Simon SD. Longitudinal study of semen
quality of unexposed workers. Sperm motility characteristics. Journal of
Andrology. 1991; 12(2):126-31. PMID: 2050580

{[}13{]} Schrader SM, Breitenstein MJ, Turner TW, Simon SD. Longitudinal
study of semen quality of unexposed workers: seminal plasma
characteristics. in Comparative spermatology 20 years after. Baccetti B,
editor. New York NY: Serno Symposia Publications; 1991.

{[}12{]} Schrader SM, Turner TW, Simon SD. Sources of variation of the
sperm penetration assay under field study conditions ARTA. 1991; 2(Suppl
4):63-74.

{[}11{]} Susten AS, Niemeier RW, Simon SD. In vivo percutaneous
absorption studies of volatile organic solvents in hairless mice. II.
Toluene, ethylbenzene and aniline. Journal of Applied Toxicology : JAT.
1990; 10(3):217-25. PMID: 2380484

{[}10{]} Ryan AS, Wysong JL, Martinez GA, Simon SD. Duration of
breast-feeding patterns established in the hospital. Influencing
factors. Results from a national survey. Clinical Pediatrics. 1990;
29(2):99-107. PMID: 2302909

{[}9{]} Schrader SM, Turner TW, Simon SD. Longitudinal study of semen
quality of unexposed workers: sperm head morphometry. Journal of
Andrology. 1990; 11(1):32-9. PMID: 2312397

{[}8{]} Simon SD, Lesage JP. Assessing the accuracy of ANOVA
calculations in statistical software. Computational Statistics and Data
Analysis. 1990; 8(3):325-32.

{[}7{]} Small LH, Simon SD, Goldberg JS. Lexical stress and lexical
access: homographs versus nonhomographs. Perception \& Psychophysics.
1988; 44(3):272-80. PMID: 3174358

{[}6{]} Schrader SM, Turner TW, Breitenstein MJ, Simon SD. Longitudinal
study of semen quality of unexposed workers. I. Study overview.
Reproductive Toxicology (Elmsford, N.Y.). 1988; 2(3-4):183-90. PMID:
2980344

{[}5{]} Simon SD, Lesage JP. The impact of collinearity involving the
intercept term on the numerical accuracy of regression. Computational
Statistics and Data Analysis. 1988; 7:197-209.

{[}4{]} Simon SD, Lesage JP. Benchmarking numerical accuracy of
statistical algorithms. Computational Statistics and Data Analysis.
1988; 7:197-209.

{[}3{]} Simon SD. How to illustrate numerical accuracy problems on any
computer. Mathematics and Computer Education. 1987; 21(1):11-5.

{[}2{]} Simon SD. To err isn't only human. Computer Language. 1986;
3(3):11-15.

{[}1{]} Lesage JP, Simon SD. Numerical accuracy of statistical
algorithms for microcomputers. Computational Statistics and Data
Analysis. 1985; 3:47-57.

\subsection{List of regional/national presentations (first author
only)}\label{list-of-regionalnational-presentations-first-author-only}

{[}81{]} Simon SD. ``What software, for what test, and why,'' session
chair at the Conference on Statistical Practice, New Orleans LA,
February 2024.

{[}80{]} Simon SD. ``Setting up an independent statistical consulting
practice,'' a poster presentation at the Conference on Statistical
Practice, New Orleans LA, February 2024.

{[}79{]} Simon SD. ``Power and sample size,'' a 90 minute guest lecture
for the UMKC Clinical Trials class, February 2024.

{[}78{]} Simon SD. ``The exterior housing conditions survey. A simple,
inexpensive, and non-intrusive approach to identify lead poisoning
risks,'' a 15 minute presentation for the National Housing Performance
Conference, Seattle WA, April 2023.

{[}77{]} Simon SD. ``What software, for what test, and why,'' session
chair at the Joint Statistical Meetings, Toronto, Canada, August 2023.

{[}76{]} Simon SD. ``Clinical statistics for non-statisticians,'' a 9
hour webinar presented for Marcus Evans, Virtual format, June 2023,
repeated December 2023 and March 2024.

{[}75{]} Simon SD. ``Exact and randomization tests,'' a 90 minute
webinar presented for The Analysis Factor, Virtual format, June 2023.

{[}74{]} Simon SD. ``Power and sample size,'' a 90 minute guest lecture
for the UMKC Clinical Trials class, February 2023.

{[}73{]} Simon SD, Wilson N. ``A summary of the Impact Lead-KC grant,''
a 30 minute presentation to the DBHI Graduate Research Seminar,
September 2022.

{[}72{]} Simon SD. ``Effective Reporting of Statistical Results by
Consultants and Collaborators,'' session chair at the Joint Statistical
Meetings, Washington DC, August 2022.

{[}71{]} Simon SD. ``A Gentle Introduction to Bootstrapping,'' a 90
minute webinar presented for The Analysis Factor, Virtual format, August
2022.

{[}70{]} Simon SD (chair). ``What I Know Now that I Wish I Knew Before I
Started a Career in Academic Consulting,'' a 90 minute panel discussion
for the Joint Statistical Meetings, in Washington DC, August 2022.

{[}69{]} Simon SD. ``Analysis of Means,'' a 90 minute webinar presented
for The Analysis Factor, Virtual format, April 2022.

{[}68{]} Simon SD. ``How to Work with Every Client--Or Not,'' a panel
discussion for the Conference on Statistical Practice, Virtual format,
February 2021.

{[}67{]} Simon SD. ``Survival analysis,'' 24 hours of webinar
presentations spread across eight weeks presented for The Analysis
Factor, Virtual format, April-June 2018. This was repeated
September-November 2018 and September-November 2020.

{[}66{]} Simon SD. ``Business Side of Setting Up an Independent
Statistical Consulting Practice,'' a 15 minute presentation for the
Joint Statistical Meetings, Virtual format, August 2020.

{[}65{]} Simon SD. ``Improving Your Scatterplots,'' a 90 minute webinar
presented for The Analysis Factor, Virtual format, April 2020.

{[}64{]} Simon SD. ``Validity and Reliability,'' a 90 minute webinar
presented for The Analysis Factor, Virtual format, November 2019.

{[}63{]} Simon SD. ``Good Enough Practices for Statistical Computing,''
a 90 minute webinar presented for The Analysis Factor, Virtual format,
August 2019.

{[}62{]} Simon SD. ``Ebenezer Scrooge, Data Scientist,'' a poster
presentation for the UMKC Faculty Research Symposium, Kansas City, MO,
April 2019.

{[}61{]} Simon SD. ``Survival analysis,'' 24 hours of webinar
presentations spread across eight weeks presented for The Analysis
Factor, April-June 2018. This was repeated in September-November 2018.

{[}60{]} Simon SD. ``Mining the Electronic Health Record,'' a 10 minute
presentation for the Midwest Bioinformatics Symposium, Kansas City, MO,
April 2019.

{[}59{]} Simon SD. ``What's the best statistical package,'' a 90 minute
webinar presented for The Analysis Factor, Virtual format, February
2018.

{[}58{]} Simon SD. ``Design and analysis of a meta-analysis,'' a 90
minute webinar presented for The Analysis Factor, Virtual format,
November 2018.

{[}57{]} Simon SD. ``Finding customers for your independent consulting
practice,'' a 90 minute webinar presented for the Statistical Consulting
Section of the American Statistical Association, Virtual format,
November 2018.

{[}56{]} Simon SD. ``Business essentials for starting an independent
consulting practice,'' a 90 minute webinar presented for the Statistical
Consulting Section of the American Statistical Association, August 2018.

{[}55{]} Simon SD. ``Tests of equivalence and non-inferiority,'' a 90
minute webinar presented for The Analysis Factor, Virtual format, April
2018.

{[}54{]} Simon SD. ``Mining the electronic health record. Why and how,''
a 30 minute presentation for the regional Frontiers in Biostatistics
conference, Kansas City, MO, April 2018.

{[}53{]} Simon SD. ``How to Write a Data Analysis Plan for a Grant
Proposal,'' a 60 presentation for the Atterbury Student Success Center,
Kansas City, MO, April 2018.

{[}52{]} Simon SD. ``Using transformations to improve your linear
regression model,'' a 90 minute webinar presented for The Analysis
Factor, Virtual format, March 2018.

{[}51{]} Simon SD. ``Setting up your independent consulting practice,''
a 90 minute talk for the Center for Research Methods and Data Analysis,
Lawrence, KS, February 2018.

{[}50{]} Simon SD. ``Good Computing Practices for Researchers,'' a 90
minute webinar presented for The Analysis Factor, November 2017, also
presented for the KUMC Biostatistics Journal Club, November 2017.

{[}49{]} Simon SD. ``Simulating clinical trials,'' an invited 20 minute
presentation for the Applied Stochastic Modeling and Data Analysis
conference, London, England, June 2017.

{[}48{]} Simon SD. ``LASSO Regression,'' a 90 minute webinar presented
for The Analysis Factor, Virtual format, November 2016.

{[}47{]} Simon SD. ``Cox Regression,'' a 90 minute webinar presented for
The Analysis Factor, Virtual format, September 2016.

{[}46{]} Simon SD. ``Formal Programs Teaching Critical Appraisal in
Evidence-Based Medicine,'' an invited 20 minute presentation for the
Joint Statistical Meeting, Chicago, IL, August 2016.

{[}45{]} Simon SD. ``Introduction to Kaplan-Meier Curves,'' a 90 minute
webinar presented for The Analysis Factor, Virtual format, April 2016.

{[}44{]} Simon SD. ``A Gentle Introduction to Bayesian Data Analysis,''
a 90 minute webinar presented for The Analysis Factor, Virtual format,
June 2015.

{[}43{]} Simon SD. ``The Hedging Hyperprior,'' a 30 minute presentation
for the regional Frontiers in Biostatistical Methods Conference, Kansas
City MO, April 2015.

{[}42{]} Simon SD. ``Setting Up and Running and Independent Statistical
Consulting Practice,'' an invited four hour short course for the
Conference on Statistical Practice, New Orleans LA, February 2015.

{[}41{]} Simon SD, Sanders S. ``Statistics for Medical Librarians,'' an
invited three hour short course for the Medical Librarians Association
national meeting, Chicago IL, May 2014.

{[}40{]} Simon SD, ``Selecting an Appropriate Sample Size'' and
``Conducting a Pilot Study,'' part of a three hour short course, Writing
a CAM Grant, for the International Research Conference on Complementary
Medicine, Miami FL, May 2014.

{[}39{]} Simon SD. ``What is a P-Value?'' an invited one hour short
course for the Midwest Society of Pediatric Research, October 2013,
Minneapolis MN.

{[}38{]} Simon SD. ``Data sources for a proposed course on secondary
data analysis,'' an invited 20 minute presentation for the Joint
Statistical Meetings, Montreal, Canada.

{[}37{]} Simon SD, Sanders S. ``Statistics for Medical Librarians,'' an
invited three hour short course for the Medical Librarians Association
national meeting, Seattle WA, May 2012.

{[}36{]} Simon SD, ``Selecting an Appropriate Sample Size'' and
``Characteristics of a Good Statistical Consultant,'' part of a three
hour short course, Writing a CAM Grant, for the International Research
Conference on Complementary Medicine, Portland OR, May 2012.

{[}35{]} Simon SD. ``Promoting Your Consulting Career,'' an invited two
hour short course for the Conference on Statistical Practice, Orlando
FL, February 2012.

{[}34{]} Simon SD. ``How Independent Consulting is Different,'' part of
a panel discussion, Statistical Consulting - Working Collaboratively
Across Disciplines, for the Conference on Statistical Practice, Orlando
FL, February 2012.

{[}33{]} Simon SD, Sanders S. ``Statistics for Medical Librarians,'' an
invited three hour short course for the Southeast Region of Medical
Librarians Association annual meeting, Augusta GA, May 2012.

{[}32{]} Simon SD. ``How Independent Consulting is Different,'' part of
a panel discussion, Successful Statistical Consulting: The
Practicalities, for the Joint Statistics Meeting, Miami Beach FL, August
2011.

{[}31{]} Simon SD. ``Using Social Media to Promote Your Consulting
Career,'' a lunchtime roundtable discussion for the Joint Statistics
Meeting, Miami Beach FL, August 2011.

{[}30{]} Simon SD. ``What is a P-Value?'' an invited one hour short
course for the Midwest Society of Pediatric Research, October 2010, Iowa
City IA.

{[}29{]} Simon SD. ``Discussion of Morphology Consensus Data,'' an
invited two hour presentation including near real-time data analysis of
classroom exercises for the Andrology Lab Workshop at the Andrology
Society national meeting, April 2010, Houston, TX.

{[}28{]} Simon SD. ``What Do All These Numbers Mean? Confidence
Intervals and P-Values,'' an invited one hour short course for the
Midwest Society of Pediatric Research, October 2009, Chicago IL.

{[}27{]} Simon SD. ``Discussion of Morphology Consensus Data,'' an
invited two hour presentation including near real-time data analysis of
classroom exercises for the Andrology Lab Workshop at the Andrology
Society national meeting, April 2009, Philadelphia, PA.

{[}26{]} Simon SD, Gajewski BJ ``Bayesian tools for planning and
monitoring accrual rates in clinical trials,'' an invited 30 minute oral
presentation for the regional Frontiers in Biostatistical Methods
Conference, April 2008, Kansas City, MO.

{[}25{]} Simon SD. ``Assessing the Safety of Pediatric Treatments,'' an
invited four hour short course for the Signal Detection and Risk
Management conference, December 2007, Amsterdam, The Netherlands.

{[}24{]} Simon SD. ``Control charts for continuous monitoring of the
number needed to harm,'' an invited 35 minute oral presentation for the
Signal Detection and Risk Management conference, December 2007,
Amsterdam, The Netherlands.

{[}23{]} Simon SD. ``Use of Diagnostic Tests for Making Clinical
Decisions,'' an invited 30 minute oral presentation for the Annual
Scientific Meeting of the American College of Allergy, Asthma \&
Immunology, November 2007, Dallas TX.

{[}22{]} Simon SD. ``Statistical Evidence: Who Was Left Out,'' an
invited 1.5 hour short course for the Midwest Society for Pediatric
Research, October 2007, Indianapolis, IN.

{[}21{]} Simon SD, Gajewski B ``Predicting Accrual in Clinical Trials:
Bayesian Posterior Predictive Distributions,'' a contributed poster for
the Joint Statistical Meetings, July/August 2007, Salt Lake City UT.

{[}20{]} Simon SD, Schrader SM. ``Using Statistics to Monitor and
Improve Quality,'' an invited three hour short course (I was responsible
for two hours) for the 2007 Andrology Lab Workshop for the American
Society for Andrology, April 2007, Tampa Bay, FL.

{[}19{]} Simon SD. ``Control charts for continuous monitoring of the
number needed to harm,'' an invited 35 minute oral presentation for the
Signal Detection and Data Mining conference, December 2006, London,
England.

{[}18{]} Simon SD. ``Signal Detection Strategies for Paediatric
Treatments,'' an invited 3 hour short course for the Signal Detection
and Data Mining conference, December 2006, London, England.

{[}17{]} Simon SD. ``Statistics Overview,'' an invited two hour
presentation for the Midwest Society for Pediatric Research, October
2006, Indianapolis IN.

{[}16{]} Simon SD. ``Statistical Evidence: Apples or Oranges?'' an
invited 1.5 hour short course for the Midwest Society for Pediatric
Research, October 2005, St.~Louis MO.

{[}15{]} Simon SD. ``What Can Alternative Medicine Teach Us About
Evidence-Based Medicine?'' an invited 1.5 hour short course for the
Midwest Society for Pediatric Research, October 2004, St.~Louis MO.

{[}14{]} Simon SD. ``Four Tips for Developing an Educational Web Site,''
a contributed poster presentation for the Joint Statistical Meetings,
August 1998, Dallas TX.

{[}13{]} Simon SD. ``Medical Statistics Case Studies on the Web,'' a
contributed 15 minute oral presentation for the Joint Statistical
Meetings, August 1997, Anaheim CA.

{[}12{]} Simon SD. ``Winning Ways of Laboratory Data Presentation:
Analyzing and Presenting Your Laboratory Data,'' Course Co-director and
Instructor for an 8 hour Andrology workshop. As an instructor, I
presented two separate talks and supervised several computer exercises.
American Society for Andrology, April 1996, Minneapolis, MN.

{[}11{]} Simon SD and many others. ``NIOSH Training Course 553:
Statistical Analysis for Industrial Hygiene Sampling and Decision
Making,'' which was offered on an annual basis from 1988 through 1993. I
delivered approximately 6 to 8 hours of lectures in a 36 hour class,
covering topics such as Normal and Lognormal Distributions;
Distributions of Means and Errors in Exposure Measurements; Confidence
and Tolerance Limits; and Selection of Workers and Times to be Sampled.

{[}10{]} Simon SD, Schrader SM. ``Statistical Methods and Experimental
Design in the Reproductive Laboratory,'' a contributed oral presentation
for the Annual Meeting of the American Society for Reproductive
Medicine, October 1995, Seattle WA.

{[}9{]} Simon SD, Schrader SM, Turner TW. ``Statistical Methods for
Human Reproductive Toxicology Research,'' an invited presentation at the
Biometric Society meeting, April 1994, Cleveland OH.

{[}8{]} Simon SD, Schrader SM, Turner TW. ``The Use of Tolerance
Intervals for Assessing Male Reproductive Potential,'' a contributed 15
minute presentation for the Joint Statistical Meetings, August 1989,
Washington, DC.

{[}7{]} Simon SD. ``SAS/Graph on the PC: User Experiences,'' a
round-table discussion section held at the SAS Users Group International
conference, April 1990, in Nashville, TN.

{[}6{]} Simon SD and many others. ``Methods for semen studies in
humans,'' held in Washington DC, October 1989. I led discussion in three
areas: ``Assessing variability between and within individuals'';
``Design considerations for parameters measured on individual cells'';
and ``Recoding continuous semen quality data as a binary response: the
application of tolerance intervals to morphometry data.''

{[}5{]} Simon SD, Lesage JP. ``Studying Numerical Accuracy as a
Controlled Experiment,'' a contributed presentation for the Conference
on the Interface Between Computer Science and Statistics, March 1987,
Philadelphia PA and published in the conference proceedings.

{[}4{]} Simon SD, Quinn LM. ``Using Some Turbo Pascal Programs to
Simulate Random Sampling in the Classroom,'' a contributed presentation
for the Winter Conference of the American Statistical Association,
January 1987, Orlando FL. There is a very similar talk by the same
authors, ``Simulating Random Sampling,'' a contributed presentation for
the Mathematics and Statistics conference, September 1985, Oxford OH.

{[}3{]} Simon SD. ``The Median Star: an Alternative to the Tukey
Resistant Line,'' a contributed 15 minute presentation for the Joint
Statistical Meetings, August 1986, Chicago IL and published in the
Proceedings of the Statistical Computing Section.

{[}2{]} Simon SD. ``Quick Estimates of Slope Based on the Cox and Stuart
Test for Trend,'' a contributed 15 minute presentation for the Joint
Statistical Meetings, August 1985, Las Vegas NV.

{[}1{]} Simon SD, Lesage JP. ``The Impact of Centering and Scaling on
Numerical Accuracy of Regression Algorithms,'' a contributed
presentation for the Applied Simulation and Modeling conference, June
1985, Montreal Canada and published in the conference proceedings.

Note: This resume was created using RMarkdown and was printed on
Thursday, March 20, 2025. With RMarkdown, you lose some of the fancy
formatting available in other programs, but the layout is clean and
simple and (most importantly) easy to maintain. This is a full resume
with complete lists of grants, publications, etc. Some of the lists are
very long. You can view an abbreviated resume at

\url{https://github.com/pmean/resume/blob/master/results/resume.pdf}

You can find the latest version of THIS resume at

\url{https://github.com/pmean/resume/blob/master/results/full-resume.pdf}

\end{document}
